\documentclass[11pt,letterpaper]{article}
\usepackage[margin=1in]{geometry}
\usepackage{graphicx} % Required for inserting images
\usepackage{tikz}
\usepackage{amsthm,amsmath,amssymb}
\usepackage{dsfont}
\usepackage{algorithm}
\usepackage{algpseudocode}
\usepackage{xcolor}
\usepackage{soul}
\usepackage{enumitem}

\usepackage[hidelinks]{hyperref}
\usepackage{cleveref}
\usepackage{etoolbox}

% Theorems
\newtheorem{theorem}{Theorem}[section]
\newtheorem{corollary}[theorem]{Corollary}
\newtheorem{lemma}[theorem]{Lemma}
\newtheorem{claim}[theorem]{Claim}
\theoremstyle{definition}
\newtheorem{definition}[theorem]{Definition}
\newtheorem{remark}[theorem]{Remark}
\newtheorem{observation}[theorem]{Observation}

\crefname{theorem}{Theorem}{Theorems}
\Crefname{theorem}{Theorem}{Theorems}
\crefname{corollary}{Corollary}{Corollaries}
\Crefname{corollary}{Corollary}{Corollaries}
\crefname{lemma}{Lemma}{Lemmas}
\Crefname{lemma}{Lemma}{Lemmas}
\crefname{claim}{Claim}{Claims}
\Crefname{claim}{Claim}{Claims}
\crefname{definition}{Definition}{Definitions}
\Crefname{definition}{Definition}{Definitions}
\crefname{remark}{Remark}{Remarks}
\Crefname{remark}{Remark}{Remarks}
\crefname{observation}{Observation}{Observations}
\Crefname{observation}{Observation}{Observations}

% The environments above share the `theorem` counter, and cleveref names a label
% after its counter -- so without these aliases every \cref to a lemma, claim,
% etc. comes out as "Theorem". \crefalias is local to the environment's group,
% so each one is typed correctly while the shared numbering is left alone.
\AtBeginEnvironment{corollary}{\crefalias{theorem}{corollary}}
\AtBeginEnvironment{lemma}{\crefalias{theorem}{lemma}}
\AtBeginEnvironment{claim}{\crefalias{theorem}{claim}}
\AtBeginEnvironment{definition}{\crefalias{theorem}{definition}}
\AtBeginEnvironment{remark}{\crefalias{theorem}{remark}}
\AtBeginEnvironment{observation}{\crefalias{theorem}{observation}}

\newcommand{\reals}{\mathbb{R}}
\newcommand{\Span}{\mathrm{span}}
\newcommand{\coreq}{\mathrm{coreq}}
\newcommand{\cone}{\mathrm{cone}}

\newcommand{\argmax}{\operatorname*{arg\,max}}
\newcommand{\argmin}{\operatorname*{arg\,min}}
\newcommand{\tr}{\operatorname{tr}}

\newcommand{\calA}{\mathcal{A}}
\newcommand{\calB}{\mathcal{B}}
\newcommand{\calC}{\mathcal{C}}
\newcommand{\calD}{\mathcal{D}}
\newcommand{\calE}{\mathcal{E}}
\newcommand{\calF}{\mathcal{F}}
\newcommand{\calG}{\mathcal{G}}
\newcommand{\calH}{\mathcal{H}}
\newcommand{\calI}{\mathcal{I}}
\newcommand{\calJ}{\mathcal{J}}
\newcommand{\calK}{\mathcal{K}}
\newcommand{\calL}{\mathcal{L}}
\newcommand{\calM}{\mathcal{M}}
\newcommand{\calN}{\mathcal{N}}
\newcommand{\calO}{\mathcal{O}}
\newcommand{\calP}{\mathcal{P}}
\newcommand{\calQ}{\mathcal{Q}}
\newcommand{\calR}{\mathcal{R}}
\newcommand{\calS}{\mathcal{S}}
\newcommand{\calT}{\mathcal{T}}
\newcommand{\calU}{\mathcal{U}}
\newcommand{\calV}{\mathcal{V}}
\newcommand{\calW}{\mathcal{W}}
\newcommand{\calX}{\mathcal{X}}
\newcommand{\calY}{\mathcal{Y}}
\newcommand{\calZ}{\mathcal{Z}}

\newcommand{\bfa}{\mathbf{a}}
\newcommand{\bfb}{\mathbf{b}}
\newcommand{\bfc}{\mathbf{c}}
\newcommand{\bfd}{\mathbf{d}}
\newcommand{\bfe}{\mathbf{e}}
\newcommand{\bff}{\mathbf{f}}
\newcommand{\bfg}{\mathbf{g}}
\newcommand{\bfh}{\mathbf{h}}
\newcommand{\bfi}{\mathbf{i}}
\newcommand{\bfj}{\mathbf{j}}
\newcommand{\bfk}{\mathbf{k}}
\newcommand{\bfl}{\mathbf{l}}
\newcommand{\bfm}{\mathbf{m}}
\newcommand{\bfn}{\mathbf{n}}
\newcommand{\bfo}{\mathbf{o}}
\newcommand{\bfp}{\mathbf{p}}
\newcommand{\bfq}{\mathbf{q}}
\newcommand{\bfr}{\mathbf{r}}
\newcommand{\bfs}{\mathbf{s}}
\newcommand{\bft}{\mathbf{t}}
\newcommand{\bfu}{\mathbf{u}}
\newcommand{\bfv}{\mathbf{v}}
\newcommand{\bfw}{\mathbf{w}}
\newcommand{\bfx}{\mathbf{x}}
\newcommand{\bfy}{\mathbf{y}}
\newcommand{\bfz}{\mathbf{z}}

\newcommand{\cost}{\mathrm{COST}}
\newcommand{\ALG}{\mathrm{ALG}}
\newcommand{\OPT}{\mathrm{OPT}}
\newcommand{\Flow}{\textsc{Flow}}
\newcommand{\ouralg}{\textsc{LineTracker}}

\newcommand{\EE}{\mathbb{E}}
\newcommand{\Var}{\mathrm{Var}}
\newcommand{\Conv}{\mathrm{Conv}}
\newcommand{\One}{\mathds{1}}

\newcommand{\interval}[1]{\lfloor #1 \rceil}
\newcommand{\dis}[1]{#1}
\newcommand{\off}{\mathrm{off}}
\newcommand{\on}{\mathrm{on}}
\newcommand{\org}{\mathrm{org}}
\newcommand{\rel}{\mathrm{rel}}
\newcommand{\new}{\mathrm{new}}

\newcommand{\solution}{\medskip\noindent\textbf{Solution.}\par\medskip}

\title{47-851: Modern Convex Optimization \\ Homework 2}
\author{Poon Tze-Yang}
\date{November 2026}


\begin{document}

\maketitle

\noindent Lecturer: Karan Singh

\begin{enumerate}[series=q, start=4]
\item \textbf{Eigenvalues and SDPs (15 points)}
\begin{enumerate}[label=\textbf{Part \Alph*.}, leftmargin=*]
\item Let $a_i \sim \calN(\mu_i, \Sigma_i)$ be multivariate normal random vectors. Recast the following problem as an SDP assuming $1/2 \leq \eta < 1$:
\[
\max_x c^T x \quad \text{subject to} \quad \Pr(a_i^T x \leq b_i) \geq \eta \quad (i = 1, \ldots, m).
\]
\item Recast the following problem as an SDP with strict inequality constraints allowed:
\[
\min_X \frac{\lambda_{\max}(X)}{\lambda_{\min}(X)} \quad \text{subject to} \quad \langle A_i, X \rangle = b_i \quad (i = 1, \ldots, m), \quad X \succ 0.
\]
\item Recast the following problem as an SDP:
\[
\min_{X \in \mathbb{S}^{n \times n}} \sum_{j=1}^k \lambda_j(X) \quad \text{subject to} \quad \langle A_i, X \rangle = b_i \quad (i = 1, \ldots, m),
\]
where $\lambda_1(X) \geq \ldots \geq \lambda_n(X)$ are the eigenvalues.

\emph{Hint:} Recast $\max_{U \in \reals^{n \times k}, U^\top U = I} \operatorname{tr}(U^\top X U) = \sum_{i=1}^k \lambda_i(X)$ as an SDP and consider the dual.
\end{enumerate}
\end{enumerate}

\solution

A.
The random variable $a_i^\top x$ is the projection of a multivariate normal vector onto $x$, so it is normally distributed.
We have that $\EE[a_i^\top x] = \mu_i^\top x$ and $\Var(a^\top_i x) = \EE[(a^\top_i x - \mu_i^\top x)^2]=\EE[x^\top (a_i - \mu_i)(a_i-\mu_i)^\top x] = x^\top \Sigma_i x$ by linearity of expectation and the fact that $\Sigma_i = \EE[(a_i - \mu_i)(a_i - \mu_i)^\top]$.

The covariance matrix $\Sigma_i$ is PSD: for any $x$, we have $x^\top \Sigma_i x = \EE[x^\top (a_i - \mu_i)(a_i-\mu_i)^\top x] = \EE[(x^\top (a_i - \mu_i))^2] \geq 0$.
Thus, we can write $\Sigma_i = L_i L_i^\top$ for some $L_i$ by the Cholesky decomposition.

Let $\Phi$ be the CDF of the standard normal distribution. Then, we have
\[
\Pr(a^\top_i x\leq b) = \Pr(\calN(\mu_i^\top x, x^\top \Sigma_i x) \leq b) = \Phi\left(\frac{b - \mu_i^\top x}{\sqrt{x^\top \Sigma_i x}}\right) = \Phi\left(\frac{b - \mu_i^\top x}{\left\|L_i^\top x\right\|_2}\right).
\]

Assume $\left\|L_i^\top x\right\|_2 > 0$, letting $\kappa_i = \Phi^{-1}(\eta_i)$, we have that the following condition is equivalent to the original inequality:
\[
    b-\mu_i^\top x \geq \left\|L_i^\top x\right\|_2 \kappa_i.
\]
Note that if $\left\|L_i^\top x\right\|_2 = 0$, then $\Pr(a^\top_i x\leq b) = \Pr(\mu^\top_i x\leq b)$, which is 1 when $\mu^\top_i x \leq b$ and 0 otherwise. 
Thus, since $\eta_i\geq \frac{1}{2}$, the original constraint is only satisfied when $b-\mu^\top_i x\geq 0$, which is equivalent to the above constraint.

Thus, we can rewrite the problem as 
\[
\max_x c^T x \quad \text{subject to} \quad \frac{b}{\kappa_i} -\frac{1}{\kappa_i}\mu_i^\top x \geq \left\|L_i^\top x\right\|_2 \quad (i = 1, \ldots, m),
\]
which is a second-order code program. Using material from the course, we can write this as an SDP.
We give the method explicitly here:

Since the norms of each constraint are non-negative, squaring each constraint is equivalent, so each constraint can be rewritten as
\[
 \left(\frac{b}{\kappa_i} -\frac{1}{\kappa_i}\mu_i^\top x\right)^2  - x^\top L_i L_i^\top x \geq 0 \quad \Leftrightarrow \quad \frac{b}{\kappa_i} -\frac{1}{\kappa_i}\mu_i^\top x - \frac{x^\top L_i L_i^\top x}{\frac{b}{\kappa_i} -\frac{1}{\kappa_i}\mu_i^\top x} \geq 0
 \]
If $\frac{b}{\kappa_i} -\frac{1}{\kappa_i}\mu_i^\top x \neq 0$, by Schur complement lemma, this is equivalent to
\[
\begin{bmatrix}
\frac{b}{\kappa_i} -\frac{1}{\kappa_i}\mu_i^\top x & x^\top L_i \\
L_i^\top x & \frac{b}{\kappa_i} -\frac{1}{\kappa_i}\mu_i^\top x
\end{bmatrix} \succeq 0.
\]
Note that if $\frac{b}{\kappa_i} -\frac{1}{\kappa_i}\mu_i^\top x = 0$, then the diagonals of the above matrix are zero.
By using the fact that PSD matrices can be written as $S_{ij} = v_i^\top v_j$ for some vectors $v_i$, we have that the above matrix is PSD if and only if each $v_i = 0$, which is equivalent to $\left\|L_i^\top x\right\|_2 = 0$, which also satisfies the constraint.

Since the above matrix is linear in $x$, we have that the problem can be written as an SDP using the dual SDP formulation.

\bigskip

B.
Consider the problem
\[
\min_{X, \alpha, t} t \quad \text{subject to} \quad \langle A_i, \alpha X \rangle = b_i \quad (i = 1, \ldots, m), \quad \alpha X - I \succeq 0, \quad t I - \alpha X \succeq 0, \quad \alpha > 0.
\]

Any solution $X,\alpha,t$ to the new problem can be mapped to a feasible solution $X'\gets \alpha X$ of the original problem.
We can verify that each constraint $\langle A_i, X' \rangle = b_i$ holds, and $X'\succ 0$ since $\alpha X-I\succeq 0$ and $\alpha>0$, so no eigenvalue of $X$ can be zero.

The minimum eigenvalue of $X$ is $\geq \alpha$, and the maximum eigenvalue of $X$ is $\leq t\alpha$. 
Hence, the minimizing $t$ is exactly the ratio between the largest and smallest eigenvalues of $X$.
Note that scaling up $X$ by $\alpha$ does not change the ratio of its eigenvalues, so the corresponding solution $X'$ in the original problem also has the same objective value.

Now suppose we have a solution $X'$ to the original problem.
We construct a feasible solution to the new problem: $\alpha\gets \lambda_{\min}(X'), t\gets \lambda_{\max}(X') \lambda_{\min}(X'), X\gets (1/\alpha)X'$.
Again, we can readily verify that the constraints $\langle A_i, \alpha X \rangle = b_i$ hold.
Clearly $\alpha>0$ since $X'\succ 0$, and $\alpha X - I \succeq 0$ and $I - \alpha X \succeq 0$ by construction of $\alpha, t$.
Thus, the constructed solution is feasible and attains the same objective value.

This suffices to show that both problems are equivalent (have the same optimum value).

% First, note that scaling $X$ by any positive constant scales all of its eigenvalues by that constant, so the ratio between the largest and smallest values does not change.
% We introduce a new variable $\alpha >0$, and require that $\alpha X\succeq 0$.
% This requires us to rewrite the equality constraints as $\langle A_i, \alpha X \rangle = \alpha b_i$.

% If we require $\alpha X-I \succeq 0$, then the minimum eigenvalue of $\alpha X$ is constrained to be at least 1.
% Similarly, for a new variable $t$, if we require $t I - \alpha X \succeq 0$, then the maximum eigenvalue of $\alpha X$ is at most $t$.

% Thus, we can rewrite the problem as
% \[
% \min_{X, \alpha, t} t \quad \text{subject to} \quad \langle A_i, \alpha X \rangle = \alpha b_i \quad (i = 1, \ldots, m), \quad \alpha X - I \succeq 0, \quad t I - \alpha X \succeq 0, \quad \alpha > 0.
% \]

% This satisfies $X\succ 0$, since the minimum eigenvalue of $\alpha X$ is at least 1 and $\alpha > 0$.
% The objective is minimized when the maximum eigenvalue of $\alpha X$ is minimized, which is equivalent to minimizing the ratio between the maximum and minimum eigenvalues of $\alpha X$, which is equal to the ratio of $X$.


\bigskip

C.
As per the hint, we consider the problem $\max_{U \in \reals^{n \times k}, U^\top U = I} \operatorname{tr}(U^\top X U) = \sum_{i=1}^k \lambda_i(X)$.
Using the cyclic property of trace, the objective is equivalently written as $\operatorname{tr}(X U U^\top)$.
Let $P = U U^\top$. 

We show that $P = U U^\top$ for $U^\top U = I$ iff $P$ has exactly $k$ eigenvalues of value 1 and $n-k$ eigenvalues of value 0.

For the forward direction, since $U^\top U=I$, we have that $P^2 = U U^\top UU^\top = UU^\top = P$. Hence, $P$ is a projection operator, so its eigenvalues are all either $0$ or $1$.
Furthermore, by the cyclic property of trace, $\tr(U U^\top) = \tr(U^\top U) = k$. Hence, $P$ has exactly $k$ eigenvalues of value $1$ and $n-k$ eigenvalues of value $0$.

For the reverse direction, let $P = Q\operatorname{diag}(I_k,0)Q^\top$ be an eigendecomposition. Taking $U=Q\begin{bmatrix}I_k\\0\end{bmatrix}$ gives $U^\top U=I_k$ and $UU^\top=P$.

Next, we encode that $P$ has exactly $k$ eigenvalues of value 1 and $n-k$ eigenvalues of value 0 with the following conditions:
(1) $P\succeq 0$ ($P$ has eigenvalues at least 0); (2) $I-P\succeq 0$ ($P$ has eigenvalues at most 1); (3) $\tr(P) = k$ (the sum of $P$'s eigenvalues is equal to $k$).
Note that these conditions need not restrict $P$ to have rank at most $k$. However, in order to maximize the objective $\sum_{i=1}^k \lambda_i(X)$, a maximizing value is achieved when $P$ projects onto the largest $k$ eigenvectors of $X$. 

This gives us the following SDP:
\[
\max_{P} \operatorname{tr}(X P) \quad \text{subject to} \quad P \succeq 0, \quad I-P \succeq 0, \quad \operatorname{tr}(P) = k.
\]
The Lagrangian is 
\begin{align*}
  L(P;W,Z,y) &= \tr(XP) + \tr(WP) + \tr(Z(I-P)) + y(k-\tr(P)) \\
  &= \tr((X+W-Z+yI)P) + \tr(Z) + yk,
\end{align*}
where $W,Z\succeq 0$.
The $\sup$ over $P$ is finite only if $X+W-Z-yI = 0$, in which case the $\sup$ is equal to $\tr(Z) - yk$. Hence, the dual SDP is
\[
\min_{y,Z}  ky + \tr(Z) \quad \text{subject to} \quad Z\succeq 0, \quad X-Z-yI \succeq 0.
\]
To show that strong duality holds, it suffices to find a $P$ that is strictly feasible for the primal SDP.
The matrix $P = \frac{k}{n}I$ is strictly feasible, since it has eigenvalues $\frac{k}{n}\in(0,1)$ and $\tr(P) = k$. Hence, strong duality holds.

Finally, we can write the original problem as
\[
\min_{y,Z,X}  ky + \tr(Z) \quad \text{subject to} \quad Z\succeq 0, \quad X-Z-yI \succeq 0, \quad \langle A_i, X \rangle = b_i \quad (i = 1, \ldots, m).  
\]

\begin{enumerate}[resume=q]
\item \textbf{Lov\'asz's bound on Shannon capacity (15 points)}

Let $\alpha(G)$ be the maximum size of an independent set. Recall
\[
\Theta(G) = \lim_{k \to \infty} \sqrt[k]{\alpha(G^k)} = \sup_{k \geq 1} \sqrt[k]{\alpha(G^k)}.
\]
\begin{enumerate}[label=\textbf{Part \Alph*.}, leftmargin=*]
\item Let $J$ be the all-ones matrix and define
\[
\nu(G) = \max_X \bigl\{\operatorname{tr}(JX) : \operatorname{tr}(X) = 1,\ X \succeq 0,\ X_{ij} = 0 \text{ if } \{i, j\} \in E\bigr\}.
\]
Show that $\nu(G) \geq \alpha(G)$.
\item Show that $\alpha(G_A \boxtimes G_B) \geq \alpha(G_A)\alpha(G_B)$. The strong product has vertex set $V_A \times V_B$. Distinct vertices $(i, k), (j, \ell)$ are adjacent exactly when
\[
(\{i, j\} \in E_A \text{ or } i = j) \quad \text{and} \quad (\{k, \ell\} \in E_B \text{ or } k = \ell).
\]
\item Show that $X, Y \succeq 0$ implies $X \otimes Y \succeq 0$. For $A \in \reals^{m \times n}$ and $B \in \reals^{p \times q}$, the Kronecker product $A \otimes B$ is defined as
\[
A \otimes B = \begin{bmatrix}
A_{11}B & A_{12}B & \cdots & A_{1n}B \\
\vdots & \vdots & \ddots & \vdots \\
A_{m1}B & A_{m2}B & \cdots & A_{mn}B
\end{bmatrix}.
\]
\item Show that $\nu(G_A \boxtimes G_B) \leq \nu(G_A)\nu(G_B)$.

\emph{Hint:} Use the dual SDP and try to \emph{Kronecker} two feasible solutions using part C.
\item Show that $\nu(G) \geq \Theta(G)$ and use this to show $\nu(C_5) = \Theta(C_5) = \sqrt{5}$.

\emph{Hint:} Numerically solving the SDP to see if it is close enough to $\sqrt{5}$ is an acceptable proof.
\end{enumerate}
\end{enumerate}

\solution
A.
Let $S\subseteq V$ be a maximum independent set of size $\alpha(G)$.
We define the matrix $X$ as follows:
\begin{align*}
    X_{ij} = \begin{cases}
        \frac{1}{\alpha(G)} &\text{if }v_i\in S \text{ and }v_j\in S,\\
        0 &\text{otherwise.}
    \end{cases}
\end{align*}
Along the diagonal, there are $\alpha(G)$ non-zero entries, so $\tr(X)=1$.
In total, there are $|S|^2$ non-zero entries, so $\tr(JX) = \alpha(G)$.
Hence, the maximizing $X$ achieves $\nu(G)\geq \alpha(G)$.

B.
Given maximum independent sets $S_A$ and $S_B$, define the set 
\[
S = \{(i,k)\in V_A\times V_B: i\in S_A\text{ and }k\in S_B\}.
\]
Immediately, $|S| = \alpha(G_A)\alpha(G_B)$.
Suppose for contradiction $(i,k)\neq(j,\ell)\in S$ and $(i,k)\sim (j,\ell)$.
Since these two vertices are distinct, either $i\neq j$ or $k\neq \ell$.
Suppose $i\neq j$, the other case is symmetric.
Then $k\sim\ell$ in $G_B$, but this is a contradiction with $k,\ell\in S_B$.
Thus, $S$ is independent in $G_A\boxtimes G_B$, so $\alpha(G_A \boxtimes G_B) \geq \alpha(G_A)\alpha(G_B)$.

C.
Let $x,y$ be any eigenvectors of $A,B$ with eigenvalues $\lambda,\mu$, respectively. 
Then 
\begin{align*}(A\otimes B) (x\otimes y) &= \begin{bmatrix}
A_{11}x_1(By) & A_{12}x_2(By) & \cdots & A_{1n}x_n(By) \\
\vdots & \vdots & \ddots & \vdots \\
A_{m1}x_1(By) & A_{m2}x_2(By) & \cdots & A_{mn}x_n(By)
\end{bmatrix} \\
&= (Ax)\otimes (\mu y) \\
&= \lambda\mu (x\otimes y).
\end{align*}
So since all eigenvalues of $A,B$ are nonnegative, all eigenvalues of $A\otimes B$ are nonnegative as well, which shows $A\otimes B\succeq 0$.

D.
For the primal problem
\[
\nu(G) = \max_X \bigl\{\operatorname{tr}(JX) : \operatorname{tr}(X) = 1,\ X \succeq 0,\ X_{ij} = 0 \text{ if } \{i, j\} \in E\bigr\}.
\]
we write the dual problem 
\[
\min_{t,Y} \{t: tI+Y-J\succeq 0, Y \text{ symmetric and } Y_{ij}=0 \text{ for }(i,j)\notin E\}.
\]
Note that $X=I/n\succ0$ is strictly feasible for the primal, so strong duality holds and the optimal dual value is $\nu(G)$.
The primal feasible set is nonempty and compact (it lies in $\{X\succeq 0 : \tr(X) = 1\}$), so the primal optimum is finite, and Slater's condition also guarantees that the dual optimum is attained.

Since $Y_{ii} = 0$ (a pair $(i,i)$ is never an edge), substituting $W = tI + Y$ is a bijection between dual feasible pairs $(t,Y)$ and symmetric $W$ with constant diagonal $t$ and zeros on non-edges. The dual becomes
\[
\nu(G) = \min_{t,W}\bigl\{t : W \succeq J,\ W_{ii} = t \text{ for all } i,\ W_{ij} = 0 \text{ for distinct non-adjacent } i,j\bigr\},
\]
with the entries of $W$ on edges unconstrained.
This form is suited to Kronecker products, which preserve zero entries and the all-ones matrix, whereas $tI+Y-J$ has entries $-1$ on non-edges.
Weak duality in this form is a one-line check: if $X$ is primal feasible and $(t,W)$ is dual feasible, then
\[
\tr(JX) \leq \tr(WX) = \sum_i W_{ii}X_{ii} + \sum_{i\neq j} W_{ij}X_{ij} = t\tr(X) = t.
\]
The inequality holds because $\tr((W-J)X) = \tr(X^{1/2}(W-J)X^{1/2}) \geq 0$. Every off-diagonal term vanishes, because either $\{i,j\}\in E$, so $X_{ij} = 0$, or $\{i,j\}\notin E$, so $W_{ij} = 0$.

Now let $(t_A, W_A)$ and $(t_B, W_B)$ be dual optimal for $G_A$ and $G_B$, so $t_A = \nu(G_A)$ and $t_B = \nu(G_B)$.
Let $W = W_A\otimes W_B$, with rows and columns indexed by $V_A\times V_B$, so that
\[
W_{(i,k),(j,\ell)} = (W_A)_{ij}(W_B)_{k\ell}.
\]
$W$ is symmetric since $(W_A\otimes W_B)^\top = W_A^\top\otimes W_B^\top$. We check that $(t_At_B, W)$ is dual feasible for $G_A\boxtimes G_B$.
\begin{itemize}
\item \emph{Diagonal.} $W_{(i,k),(i,k)} = (W_A)_{ii}(W_B)_{kk} = t_At_B$.
\item \emph{Non-edges.} Let $(i,k)\neq(j,\ell)$ be non-adjacent in $G_A\boxtimes G_B$. Negating the adjacency condition from Part B, either $i\neq j$ and $\{i,j\}\notin E_A$, or $k\neq\ell$ and $\{k,\ell\}\notin E_B$. In the first case $(W_A)_{ij} = 0$, and in the second $(W_B)_{k\ell} = 0$. Either way $W_{(i,k),(j,\ell)} = 0$.
\item \emph{$W\succeq J$.} The all-ones matrix on $V_A\times V_B$ is $J = J_A\otimes J_B$, so by bilinearity of $\otimes$,
\[
W - J = W_A\otimes W_B - J_A\otimes J_B = (W_A - J_A)\otimes W_B + J_A\otimes(W_B - J_B).
\]
By feasibility $W_A - J_A\succeq 0$ and $W_B - J_B\succeq 0$. Also $J_A = \One\One^\top\succeq 0$, and so $W_B = (W_B - J_B) + J_B\succeq 0$. By Part C both terms are PSD, hence so is their sum.
\end{itemize}
Therefore $(t_At_B, W)$ is dual feasible for $G_A\boxtimes G_B$, and weak duality gives
\[
\nu(G_A\boxtimes G_B) \leq t_At_B = \nu(G_A)\nu(G_B).
\]

E.
Identify $V^k$ with $V^{k-1}\times V$. Then $G^k = G^{k-1}\boxtimes G$, since in both graphs distinct tuples are adjacent exactly when every coordinate is equal or adjacent.
By Part D and induction on $k$, $\nu(G^k)\leq\nu(G^{k-1})\nu(G)\leq\nu(G)^k$.
Combining this with Part A applied to $G^k$,
\[
\sqrt[k]{\alpha(G^k)} \leq \sqrt[k]{\nu(G^k)} \leq \nu(G) \qquad\text{for every } k\geq 1,
\]
and taking the supremum over $k$ gives $\Theta(G)\leq\nu(G)$.

Now let $G = C_5$, with vertices $0,\dots,4$ and $i\sim i\pm 1 \pmod 5$.

\emph{Lower bound.} The five vertices $(i, 2i \bmod 5)$, $i = 0,\dots,4$, form an independent set in $C_5^2$.
To see this, take two of them with $i\neq j$ and let $d = j - i \bmod 5\in\{1,2,3,4\}$.
If $d\in\{2,3\}$, the first coordinates are distinct and non-adjacent.
If $d\in\{1,4\}$, the second coordinates differ by $2d\equiv 2$ or $3\pmod 5$, so they are distinct and non-adjacent.
Either way the two vertices are not adjacent in $C_5^2$. So $\alpha(C_5^2)\geq 5$, and
\[
\sqrt5 \leq \sqrt{\alpha(C_5^2)} \leq \Theta(C_5) \leq \nu(C_5).
\]

\emph{Upper bound.} By weak duality it suffices to exhibit one dual feasible $(t,W)$ with $t = \sqrt5$.
The rotation $i\mapsto i+1$ is an automorphism of $C_5$. Conjugating a dual feasible $W$ by its permutation matrix therefore keeps it feasible with the same $t$, and averaging over the five rotations (the feasible set is convex) gives a feasible circulant.
So we look for $W = tI + yA$, where $A$ is the adjacency matrix of $C_5$.
This $W$ has diagonal $t$ and is zero on non-adjacent pairs, so it remains to choose $t,y$ with $W - J\succeq 0$.

Let $\omega = e^{2\pi\mathrm{i}/5}$ and $v_m = (\omega^{jm})_{j=0}^{4}$ for $m = 0,\dots,4$. These vectors form a basis of $\mathbb{C}^5$ and are eigenvectors of $I$, $A$ and $J$ simultaneously:
\[
(Av_m)_j = \omega^{(j+1)m} + \omega^{(j-1)m} = (\omega^m + \omega^{-m})(v_m)_j = 2\cos(2\pi m/5)\,(v_m)_j,
\]
and $Jv_m = \bigl(\sum_{j}\omega^{jm}\bigr)\One$, which is $5v_0$ for $m = 0$ and $0$ otherwise.
To evaluate the cosines, let $u = 2\cos(2\pi/5) = \omega + \omega^{-1}$. Then $\omega^2 + \omega^{-2} = u^2 - 2$, and $1 + \omega + \omega^2 + \omega^3 + \omega^4 = 0$ gives $u + (u^2 - 2) = -1$, i.e.\ $u^2 + u - 1 = 0$. Since $u > 0$, we get
\[
2\cos(2\pi/5) = u = \tfrac{\sqrt5 - 1}{2}, \qquad 2\cos(4\pi/5) = u^2 - 2 = -1 - u = -\tfrac{\sqrt5 + 1}{2}.
\]
Hence the eigenvalues of $W - J$ are
\[
t + 2y - 5 \ \ (m = 0), \qquad t + \tfrac{\sqrt5 - 1}{2}\,y \ \ (m = 1,4), \qquad t - \tfrac{\sqrt5 + 1}{2}\,y \ \ (m = 2,3).
\]
Minimizing $t$ subject to these three being nonnegative is a two-variable LP, whose optimum makes the $m = 0$ and $m = 2,3$ constraints tight.
This gives $t = \sqrt5$, $y = \frac{5 - \sqrt5}{2}$. We now verify that this point is feasible:
\begin{align*}
t + 2y - 5 &= \sqrt5 + (5 - \sqrt5) - 5 = 0,\\
t + \tfrac{\sqrt5 - 1}{2}\,y &= \sqrt5 + \tfrac{(5 - \sqrt5)(\sqrt5 - 1)}{4} = \sqrt5 + \tfrac{6\sqrt5 - 10}{4} = \tfrac{5\sqrt5 - 5}{2} > 0,\\
t - \tfrac{\sqrt5 + 1}{2}\,y &= \sqrt5 - \tfrac{(5 - \sqrt5)(\sqrt5 + 1)}{4} = \sqrt5 - \tfrac{4\sqrt5}{4} = 0.
\end{align*}
So $W = \sqrt5\, I + \frac{5 - \sqrt5}{2}A$ satisfies $W\succeq J$, and $(\sqrt5, W)$ is dual feasible for $C_5$. Thus $\nu(C_5)\leq\sqrt5$, and together with the lower bound,
\[
\nu(C_5) = \Theta(C_5) = \sqrt5.
\]





% \begin{enumerate}[resume=q]
% \item \textbf{Grads and Dists (10 points)}
% \begin{enumerate}[label=\textbf{Part \Alph*.}, leftmargin=*]
% \item Consider a function $f : \reals^d \to \reals$ that is $\mu$-strongly convex and $L$-smooth. Let $x^* \in \argmin_x f(x)$. Prove for all $x \in \reals^d$ that
% \[
% \frac{1}{2L}\|\nabla f(x)\|^2 \leq f(x) - f(x^*) \leq \frac{1}{2\mu}\|\nabla f(x)\|^2.
% \]
% \item Prove that $f$ is $\mu$-strongly convex if and only if for all $x, y \in \reals^d$, we have
% \[
% (\nabla f(x) - \nabla f(y))^\top (x - y) \geq \mu\|x - y\|^2.
% \]
% \end{enumerate}
% \end{enumerate}

% \solution
% A.
% By definition, $f$ is $\mu$-strongly convex iff
% \[
% f(y) \geq f(x) + \nabla f(x)^\top (y-x) + \frac{\mu}{2}\|y-x\|^2.
% \]
% Rearranging:
% \[
% f(y) -f(x) \geq \nabla f(x)^\top (y-x) + \frac{\mu}{2}\|y-x\|^2.
% \]
% Taking the minimum of both sides w.r.t. $y$:
% \[
% \min_y\{f(y) -f(x)\} \geq \min_y\{\nabla f(x)^\top (y-x) + \frac{\mu}{2}\|y-x\|^2\}.
% \]
% The left hand side is equal to $f(x^*)-f(x)$.
% The right hand side is minimized when the derivative with respect to $y$ is equal to 0:
% \begin{align*}
%     \nabla f(x) + \frac{\mu}{2}(2y-2x) = 0\quad\Rightarrow\quad y -x = -\frac{1}{\mu}\nabla f(x).
% \end{align*}
% Substituting back in:
% \[
% \min_y\{\nabla f(x)^\top (y-x) + \frac{\mu}{2}\|y-x\|^2\} = -\frac{1}{2\mu}\|\nabla f(x)\| ^2.
% \]
% Combining both sides:
% \[
% f(x^*) - f(x) \geq -\frac{1}{2\mu}\|\nabla f(x)\| ^2 \quad\Rightarrow\quad f(x) - f(x^*)  \leq \frac{1}{2\mu}\|\nabla f(x)\| ^2.
% \]

% A function $f$ is $L$-smooth if
% \[
% f(y)\leq f(x) + \nabla f(x)^\top(y-x) + \frac{L}{2}\left\|(y-x)\right\|^2
% \]
% Essentially the same argument as above gives 
% \[
% f(x)-f(x^*) \geq \frac{1}{2L}\|\nabla f(x)\|^2.
% \]
% \medskip
% B.
% For the forward direction, use $\mu$-strong convexity at $x,y$:
% \begin{align*}
%     f(y) - f(x) \geq \nabla f(x)^\top (y-x) + \frac{\mu}{2}\|y-x\|^2. \\
%     f(x) - f(y) \geq \nabla f(y)^\top (x-y) + \frac{\mu}{2}\|y-x\|^2.
% \end{align*}
% Combining the two inequalities:
% \begin{align*}
%     \nabla f(y)^\top (y-x) - \frac{\mu}{2}\|y-x\|^2 \geq \nabla f(x)^\top (y-x) + \frac{\mu}{2}\|y-x\|^2 \\
%     \Rightarrow (\nabla f(x) - \nabla f(y))^\top (x - y) \geq \mu\|x - y\|^2.
% \end{align*}

% For the reverse direction:


\end{document}
